\documentclass[conference]{IEEEtran}
\IEEEoverridecommandlockouts
\usepackage[utf8]{inputenc}
\usepackage[T1]{fontenc}
\usepackage{cite}
\usepackage{amsmath,amssymb,amsfonts}
\usepackage{graphicx}
\usepackage{textcomp}
\usepackage{xcolor}
\usepackage{hyperref}
\usepackage{listings}
\usepackage{booktabs}

\lstset{
  basicstyle=\footnotesize\ttfamily,
  breaklines=true,
  frame=single,
  keywordstyle=\color{blue},
  commentstyle=\color{green!60!black},
}

\begin{document}

\title{Dijital Gardrop: Yapay Zeka Destekli Moda Sosyal Medya Platformu\\
\thanks{YZTA Bootcamp 2026 -- Takim 35 Proje Raporu}}

\author{
\IEEEauthorblockN{Ahmet Colak}
\IEEEauthorblockA{\textit{YZTA Bootcamp 2026} \\
Takim 35}
}

\maketitle

\begin{abstract}
Bu calisma, kullanicilarin dijital gardiroplarini yonetmelerine, yapay zeka destekli kombin onerileri almalarına ve moda icerikleri paylasmalarina olanak taniyan \textit{Dijital Gardrop} platformunun tasarimini, gelistirilmesini ve test surecini sunmaktadir. Sistem; Flutter tabanli cok platformlu mobil istemci, FastAPI tabanli RESTful arka uc ve yerel olarak calisan Ollama/LLaMA 3.2 dil modeli uclusunden olusmaktadir. Gelistirme surecinde karsilasilan kritik sorunlar (iOS kamera izinleri, SQLite kisit ihlalleri, oturum kaliciligi eksikligi, API yonlendirme hatalari) sistematik bicimde cozulmustur.
\end{abstract}

\begin{IEEEkeywords}
Flutter, FastAPI, SQLite, Ollama, LLaMA, Moda AI, Sosyal Medya, Mobil Uygulama
\end{IEEEkeywords}

\section{Giris}

Moda endustrisi, dijital donusumun en hizli yasandigi sektorler arasinda yer almaktadir. Kullanicilarin dijital ortamda kiyafet kombinleri olusturmalari, baskalariyla paylasmalar ve yapay zeka destekli stilist hizmetlerinden yararlanmalari giderek daha yaygin bir ihtiyac haline gelmektedir.

\textit{Dijital Gardrop} projesi; bir mobil uygulama olarak gelistirilen, kiyafet yonetimi, sosyal medya ve yapay zeka onerilerini tek cati altinda birlestiren bir platform sunmaktadir. Proje, YZTA Bootcamp 2026 kapsaminda gelistirilmis olup bu rapor, teknik mimariden karsilasilan sorunlara ve uygulanan cozumlere kadar tum gelistirme surecini kapsamaktadir.

\section{Sistem Mimarisi}

\subsection{Genel Bakis}

Sistem, birbirine RESTful API araciligiyla baglanan uc ana katmandan olusmaktadir:

\begin{enumerate}
  \item \textbf{Istemci Katmani:} Flutter (Dart) ile gelistirilen iOS/Android uygulamasi.
  \item \textbf{Uygulama Katmani:} Python 3.9+ ve FastAPI ile gelistirilen REST arka uc.
  \item \textbf{AI Katmani:} Yerel agda calisan Ollama sunucusu uzerindeki LLaMA 3.2 modeli.
\end{enumerate}

\subsection{Kullanilan Teknolojiler}

\begin{table}[h]
\caption{Kullanilan Teknoloji Yigini}
\label{tab:tech}
\centering
\begin{tabular}{lll}
\toprule
\textbf{Katman} & \textbf{Teknoloji} & \textbf{Surum} \\
\midrule
Mobil Istemci & Flutter & 3.10+ \\
Dil (Istemci) & Dart & 3.0+ \\
Durum Yonetimi & Riverpod & 2.4.9 \\
Gorsel Onbellek & cached\_network\_image & 3.3+ \\
Resim Secici & image\_picker & 1.2+ \\
Oturum Kaliciligi & shared\_preferences & 2.2+ \\
\midrule
REST Arka Uc & FastAPI & 0.110+ \\
Dil (Arka Uc) & Python & 3.9+ \\
Veritabani & SQLite & 3.x \\
ASGI Sunucusu & Uvicorn & 0.29+ \\
HTTP Istemcisi & httpx & 0.27+ \\
Statik Dosyalar & FastAPI StaticFiles & -- \\
\midrule
Yapay Zeka & Ollama (LLaMA 3.2) & 3.2 \\
\bottomrule
\end{tabular}
\end{table}

\section{Veri Modeli}

Uygulama SQLite veritabani uzerinde su tablolarla calismaktadir:

\begin{itemize}
  \item \texttt{users} -- Kullanici profilleri (user\_id, email, username, avatar, bio, followers\_count, following\_count, is\_private)
  \item \texttt{posts} -- Gonderiler (post\_id, user\_id, image\_url, caption, visibility, likes\_count, comments\_count)
  \item \texttt{post\_outfit\_items} -- Gonderi--kiyafet iliskisi (category CHECK kisiti: ust giyim, alt giyim, ayakkabi, aksesuar, dis giyim, diger)
  \item \texttt{kiyafetler} -- Dijital gardirop kiyafetleri (tur, renk, beden, marka, mevsim, temiz, foto\_url)
  \item \texttt{likes} -- Begeniler (post\_id, user\_id)
  \item \texttt{follows} -- Takip iliskisi (follower\_id, following\_id)
  \item \texttt{sohbetler} -- AI stilist sohbet gecmisi (user\_id, rol, mesaj, olusturma\_zamani)
  \item \texttt{kombin\_oneriler} -- Kombin onerileri (user\_id, baglam\_json, aciklama)
  \item \texttt{comments} -- Yorumlar (comment\_id, post\_id, user\_id, text, created\_at)
\end{itemize}

\section{API Uc Noktalari}

\begin{table}[h]
\caption{REST API Endpoint Ozeti}
\label{tab:api}
\centering
\begin{tabular}{lll}
\toprule
\textbf{Grup} & \textbf{Endpoint} & \textbf{Aciklama} \\
\midrule
Auth & POST /auth/login & Giris \\
     & POST /auth/register & Kayit \\
\midrule
Posts & POST /posts & Gonderi olustur \\
      & DELETE /posts/\{id\} & Gonderi sil \\
      & GET /posts/users/\{uid\}/posts & Kullanici gonderileri \\
\midrule
Likes & POST /posts/\{id\}/like & Begen \\
      & DELETE /posts/\{id\}/like & Begeni kaldir \\
      & POST /posts/\{id\}/comments & Yorum ekle \\
      & GET /posts/\{id\}/comments & Yorumlari listele \\
\midrule
Feed & GET /feed & Akis \\
\midrule
Users & GET /users/\{uid\} & Profil \\
      & PUT /users/me & Profil guncelle \\
      & PUT /users/me/privacy & Gizlilik ayari \\
\midrule
Follows & POST /follow & Takip et \\
        & DELETE /follow & Takipten cik \\
\midrule
Wardrobe & POST /wardrobe/items & Kiyafet ekle \\
         & GET /wardrobe/items/\{uid\} & Kiyafetleri listele \\
         & POST /wardrobe/chat & AI stilist chat \\
         & GET /wardrobe/chat/history/\{uid\} & Sohbet gecmisi \\
         & POST /wardrobe/outfit/suggest & Kombin oner \\
\midrule
Captions & POST /captions/suggest & AI aciklama oner \\
         & POST /captions/upload & Resim yukle \\
\midrule
Static & GET /static/uploads/\{file\} & Resim sun \\
\bottomrule
\end{tabular}
\end{table}

\section{Yapay Zeka Entegrasyonu}

\subsection{Ollama / LLaMA 3.2}

Baslangicta Google Gemini API kullanimi planlanmis, ancak kullanicinin yerel dil modeli tercih etmesi uzerine sistem Ollama uzerindeki LLaMA 3.2 modeline gecirilmistir. Bu gecis su avantajlari saglamistir:

\begin{itemize}
  \item Veri gizliligi -- istemler harici sunuculara gonderilmez.
  \item Sifir API maliyeti -- harici ucretlendirme yoktur.
  \item Cevrimdisi calisma -- internet baglantisi gerektirmez.
\end{itemize}

\subsection{AI Stilist Chatbot}

AI stilist bileşeni \texttt{gemini\_client.py} modulunde uygulanmistir. Model, Turkce sistem promptlari ile yonlendirilmekte; JSON cikti ayristirma basarisiz oldugunda guvenli bir geri donus mekanizmasi devreye girmektedir. Iki temel islevsellik sunulmaktadir:

\begin{lstlisting}
def sohbet_yaniti_al(gecmis, yeni_mesaj):
    # Sistem promptu + gecmis + yeni mesaj
    # JSON yanit: asistan_mesaji, baglam, hazir_mi
    ...

def kombin_onerisi_uret(baglam, temiz_kiyafetler):
    # Baglam + kiyafet listesi
    # JSON yanit: secilen_idler, aciklama
    ...
\end{lstlisting}

\section{Karsilasilan Sorunlar ve Cozumler}

\subsection{iOS Kamera Cokmesi}

\textbf{Sorun:} Kiyafet ekleme ekraninda kamera butonuna basildiginda uygulama iOS'ta cokuyordu.

\textbf{Kok Neden:} \texttt{Info.plist} dosyasinda \texttt{NSCameraUsageDescription} ve \texttt{NSPhotoLibraryUsageDescription} anahtarlari eksikti.

\textbf{Cozum:} Dort izin anahtari \texttt{Info.plist}'e eklendi: \texttt{NSCameraUsageDescription}, \texttt{NSPhotoLibraryUsageDescription}, \texttt{NSPhotoLibraryAddUsageDescription}, \texttt{NSMicrophoneUsageDescription}.

\subsection{SQLite CHECK Kisit Ihlali}

\textbf{Sorun:} Gonderi olusturulurken 500 hatasi: \textit{"Check constraint failed IN ('ust giyim', 'alt giyim', ...)"}.

\textbf{Kok Neden:} \texttt{post\_outfit\_items} tablosuna \texttt{category = 'unknown'} degeri insert ediliyordu.

\textbf{Cozum:} \texttt{post\_repository.py} dosyasinda sabit deger \texttt{'unknown'} $\rightarrow$ \texttt{'diger'} olarak guncellendi.

\subsection{Oturum Kaliciligi Eksikligi}

\textbf{Sorun:} Uygulama her acildiginda giris ekrani gosteriliyor, kullanici oturumu hatirlanmiyordu.

\textbf{Kok Neden:} \texttt{AuthProvider} sinifi yalnizca bellekte userId tutuyor, kalici depolamaya yazilmiyordu.

\textbf{Cozum:} \texttt{shared\_preferences} paketi entegre edildi. \texttt{\_AuthGate} widget'i oturum durumuna gore yonlendirme yapacak sekilde \texttt{main.dart}'a eklendi.

\subsection{Resim Yukleme Akisi}

\textbf{Sorun:} Yerel dosya yolu backend'e gonderildiginden resimler diger cihazlarda goruntulenemiyor.

\textbf{Cozum:}
\begin{enumerate}
  \item Backend'e \texttt{POST /captions/upload} multipart endpoint'i eklendi.
  \item Dosya UUID adiyla \texttt{static/uploads/} dizinine kaydedildi.
  \item Frontend resmi once yukluyor, donen URL'yi post/kiyafet kaydina yaziyor.
\end{enumerate}

\subsection{Likes/Comments 404 Hatasi}

\textbf{Sorun:} Begeni ve yorum endpoint'leri 404 donduruyordu.

\textbf{Kok Neden:} \texttt{likes.py} router'i \texttt{main.py} dosyasina hic dahil edilmemisti.

\textbf{Cozum:} \texttt{main.py}'e router kayit edildi. \texttt{POST/GET /posts/\{id\}/comments} endpoint'leri eklendi. DELETE begeni query param ile \texttt{user\_id} alacak sekilde guncellendi.

\subsection{AI Stilist Yanit Sorunu}

\textbf{Sorun:} AI Stilist ekrani "I could not generate response." mesaji gosteriyordu.

\textbf{Kok Neden 1:} Frontend \texttt{response['reply']} anahtarini okurken backend \texttt{asistan\_mesaji} anahtarini donduruyordu.

\textbf{Kok Neden 2:} Ekran \texttt{user\_id = 'user\_123'} seklinde sabit deger kullaniyordu.

\textbf{Cozum:} \texttt{ai\_stylist\_screen.dart} tamamen yeniden yazildi; dogru JSON anahtarlari, gercek kullanici ID'si ve gecmis yukleme ozelligi eklendi.

\subsection{Python 3.9 Uyumluluk Sorunlari}

\textbf{Sorun:} Backend Python 3.10+ tip soz dizimi (\texttt{list[X]}, \texttt{str | None}) kullandığından Python 3.9 ortaminda hata veriyordu.

\textbf{Cozum:} Tum tip ek aciklamalari \texttt{from typing import List, Optional} kullanilarak guncellendi.

\subsection{Gonderi Silme Sorunu}

\textbf{Sorun:} Profil ekranindaki Sil butonu yalnizca SnackBar gosteriyor, API cagrisi yapmiyordu.

\textbf{Cozum:} \texttt{DELETE /posts/\{post\_id\}?user\_id=\{uid\}} endpoint'i backend'e eklendi; \texttt{profile\_provider.dart}'a \texttt{deletePost} metodu eklendi.

\section{Ag Baglantisi ve Dagitim}

Uygulama gelistirme asamasinda yerel ag uzerinde test edilmistir. iPhone'un kisisel hotspot'u uzerinden Mac'in IP adresi (\texttt{172.20.10.13:8000}) kullanilarak baglanti saglanmistir.

\section{Sprint Ozeti}

\begin{table}[h]
\caption{Sprint Gorev Durumlari}
\label{tab:sprint}
\centering
\begin{tabular}{lll}
\toprule
\textbf{Gorev} & \textbf{Durum} & \textbf{Sprint} \\
\midrule
Proje yapisi kurulumu & Tamamlandi & Sprint 1 \\
FastAPI arka uc iskeleti & Tamamlandi & Sprint 1 \\
Flutter temel ekranlar & Tamamlandi & Sprint 1 \\
Veritabani semasi & Tamamlandi & Sprint 1 \\
\midrule
Kullanici kimlik dogrulama & Tamamlandi & Sprint 2 \\
Gonderi olusturma & Tamamlandi & Sprint 2 \\
Feed akisi & Tamamlandi & Sprint 2 \\
Dijital gardirop CRUD & Tamamlandi & Sprint 2 \\
Cok dilli destek (TR/EN) & Tamamlandi & Sprint 2 \\
\midrule
AI Stilist chat & Tamamlandi & Sprint 3 \\
Kombin onerisi & Tamamlandi & Sprint 3 \\
Resim yukleme & Tamamlandi & Sprint 3 \\
iOS kamera izinleri & Tamamlandi & Sprint 3 \\
Oturum kaliciligi & Tamamlandi & Sprint 3 \\
\midrule
Begeni sistemi & Tamamlandi & Sprint 4 \\
Yorum sistemi & Tamamlandi & Sprint 4 \\
Takip sistemi & Tamamlandi & Sprint 4 \\
Gonderi silme & Tamamlandi & Sprint 4 \\
Sohbet gecmisi & Tamamlandi & Sprint 4 \\
Profil gizlilik & Tamamlandi & Sprint 4 \\
\midrule
Uretim dagitimi & Beklemede & Sprint 5 \\
Push bildirimler & Beklemede & Sprint 5 \\
Performans optimizasyonu & Beklemede & Sprint 5 \\
\bottomrule
\end{tabular}
\end{table}

\section{Sonuc}

Bu calismada, yapay zeka destekli moda sosyal medya platformu \textit{Dijital Gardrop}'un tasarimi ve gelistirilmesi sunulmustur. Flutter, FastAPI ve Ollama/LLaMA 3.2 teknolojilerini bir araya getiren platform; kiyafet yonetimi, kombin onerisi, sosyal etkilesim ve icerik paylasimini entegre etmektedir.

Gelistirme surecinde karsilasilan teknik engeller sistematik olarak asilmistir. Sistemin yerel yapay zeka modeliyle calismasi, veri gizliligi acisindan onemli bir avantaj sunmaktadir.

\begin{thebibliography}{00}
\bibitem{flutter} Flutter Team, ``Flutter -- Build apps for any screen,'' Google LLC, 2024.
\bibitem{fastapi} S. Ramirez, ``FastAPI,'' 2024. Available: https://fastapi.tiangolo.com
\bibitem{ollama} Ollama Team, ``Ollama,'' 2024. Available: https://ollama.com
\bibitem{riverpod} R. Rousselet, ``Riverpod,'' 2024. Available: https://riverpod.dev
\bibitem{sqlite} D. R. Hipp, ``SQLite,'' 2024. Available: https://sqlite.org
\bibitem{llama3} Meta AI, ``Llama 3,'' Meta Platforms, 2024.
\end{thebibliography}

\end{document}
